Researchers \must describe in the \paper the full architecture of LLM-based tools, from standalone uses to agentic systems, including the LLM's role and its interactions with other components. Hosting and access \should be described; for time-sensitive measurements, researchers \must clarify whether local infrastructure or cloud services were used. Researchers \must publish all prompts as \supplementarymaterial, including prompt templates with representative instances and the dynamic generation process where applicable, and \must include representative examples in the \paper; sensitive content \must be anonymized. If full prompt disclosure is not feasible, summaries or examples \should be provided. The prompting strategy and prompt reuse across models and configurations \must be specified, and how the prompts were developed \should be described; for few-shot prompts, how the examples were selected \must be explained in the \paper. Researchers \must describe the configuration mechanisms used (e.g., context files, skills, subagents) and summarize in the \paper which tools were exposed; all configuration artifacts and the complete tool catalog (schemas, definitions, Model Context Protocol (MCP) servers) \should be published as \supplementarymaterial. For agentic systems, researchers \must specify the agents' roles, reasoning frameworks, and communication flows; where external tools are used, the model's reasoning, tool calls, and interactions with users or the environment \must be reported separately. For retrieval-augmented generation (RAG) and similar methods, researchers \must describe how external data was retrieved, stored, and integrated. Where legally possible, the implementation \should be open-sourced; non-disclosed proprietary components \must be acknowledged as reproducibility limitations.
